\section{Autoboxing}
\label{sec:autoboxing}

\lstinputlisting[language=Java]{generated/source/Autoboxing.java}

\inputbytecode{Autoboxing}

\paragraph{Lendo a listagem.}
\begin{itemize}
\item Em \texttt{main}, cada \texttt{Integer a = 127} virou
  \texttt{bipush 127} seguido de \texttt{invokestatic Integer.valueOf}.
  O autoboxing é uma chamada de método que o compilador escreveu por
  você.
\item Em \texttt{iguais}, \texttt{a == b} virou \texttt{if\_acmpne}:
  comparação de \emph{referências} (o prefixo \texttt{a}), não de
  valores. Para comparar valores, o compilador teria gerado
  \texttt{intValue} e \texttt{if\_icmpne}.
\end{itemize}

\paragraph{O bug que isso explica.}
O programa imprime \texttt{true} e depois \texttt{false}. A JLS (§5.1.7)
garante que o boxing de uma expressão constante entre $-128$ e $127$
devolve sempre o mesmo objeto; fora dessa faixa, não há garantia, e na
prática 128 gera dois objetos diferentes. O teste passa com valores
pequenos e falha em produção com IDs maiores. Compare wrappers com
\texttt{equals} (ou \texttt{Objects.equals}), nunca com \texttt{==}.

O caminho inverso tem outro risco. Em \texttt{lambda\$dobro\$0}
(seção~\ref{sec:invokedynamic}), \texttt{x * 2} começa com
\texttt{invokevirtual Integer.intValue}. Se \texttt{x} for \texttt{null},
essa chamada lança \texttt{NullPointerException} numa linha que não tem
nenhuma chamada de método visível no código-fonte.

\paragraph{A decisão que isso muda.}
Os itens 6 e 61 do \citetitle{livro:bloch2017} (\enquote{Avoid creating
unnecessary objects} e \enquote{Prefer primitive types to boxed
primitives}) ganham uma explicação concreta: cada \texttt{valueOf} pode
alocar. Em código quente, prefira \texttt{int}, \texttt{IntStream} e
coleções de primitivos; deixe wrappers para onde \texttt{null} tem
significado (campo opcional de um JSON, coluna anulável).
